\section{Experiments}
\label{sec:experiments}

We instantiate the framework on exoplanet atmospheres---specifically,
atmospheric composition interpretation in transmission spectra,
cloud/haze degeneracies, and instrument systematics. All numbers below
are produced by the committed pipeline under the frozen manifest
\texttt{exoplanet\_atmospheres\_v1} (cutoff 2020-12-31).

\subsection{Corpus construction}
\label{sec:experiments:corpus}

Table~\ref{tab:corpus} summarizes \corpusA{}. Literature metadata came
from NASA ADS (primary) and arXiv (supplement) under four topic queries;
cross-source deduplication by arXiv identifier removed 404 duplicate
records. A core set of 500 papers was selected by greedy,
year-stratified optimization over internal citation in-degree, keyword
coverage, and target diversity (83--84 papers per year, 2015--2020).
Object mentions were linked to the NASA Exoplanet Archive by
boundary-anchored name matching tolerant of space/hyphen variants;
naive substring matching produced systematic false positives
(``WASP-1'' firing inside ``WASP-12b'') that inflated links by
$2.3\times$ before the fix.

\begin{table}[t]
  \centering
  \caption{\corpusA{} at a glance (manifest
  \texttt{exoplanet\_atmospheres\_v1}, cutoff 2020-12-31).}
  \label{tab:corpus}
  \begin{tabular}{lr}
    \toprule
    Literature metadata (deduplicated) & 2{,}512 papers \\
    \quad from NASA ADS / arXiv & 1{,}639 / 1{,}277 \\
    Core full-text set & 500 papers \\
    Citation edges & 51{,}616 \\
    Catalog objects (NASA Exoplanet Archive) & 6{,}324 planets \\
    Observation metadata (MAST, JWST+HST spectra) & 306{,}729 records \\
    Paper--object links & 2{,}364 \\
    Catalog planets with archival spectra & 553 \\
    \bottomrule
  \end{tabular}
\end{table}

\subsection{Claim extraction: calibration before scale}
\label{sec:experiments:claims}

A domain reviewer first annotated 20 claim-dense core papers by hand,
yielding 37 gold claims with source locations, assumptions, and
uncertainties. Model extraction (from abstracts in this version) was
then calibrated against the gold set over three prompt revisions:
the initial prompt over-split coupled measurements, occasionally copied
first-person phrasing, and preferred easily-copied quantitative results
over interpretive conclusions (e.g., it captured ``mostly subsolar
H$_2$O'' across ten hot Jupiters \citep{pinhas2019} but dropped the
paper's load-bearing interpretation of subsolar oxygen or supersolar
C/O). The final prompt enforces 2--5 load-bearing claims with a priority
order (conclusion $>$ interpretation $>$ methodological challenge $>$
decisive quantity $>$ secondary measurement), merging of coupled
measurements, normalized third person, and an omission self-check.
Under it, all reviewer-cited misses were recovered, first-person copying
fell to zero, and no paper exceeded five claims (mean 3.68 across
353 claims from 100 papers; 4 methods papers correctly yielded none).

Semantic deduplication against the gold set (same paper, compatible
objects, same conclusion polarity, embedding cosine $\ge 0.62$) matched
all 37 gold claims to a model counterpart (similarity 0.63--0.97, median
0.76), producing a canonical table of 353 claims: 37
\texttt{model\_matched\_gold} and 316 \texttt{model\_only}.

\subsection{Tension detection and adjudication}
\label{sec:experiments:tension}

The high-recall candidate rule produced 161 cross-paper tension
candidates over the canonical claims (807-node, 1{,}237-edge graph).
The reviewer adjudicated the top-ranked candidates; the resulting
16 typed verdicts are themselves a finding
(Table~\ref{tab:adjudication}): only \emph{one} candidate is a genuine
observational contradiction---the HD~189733~b wind-velocity discrepancy,
${\sim}8$~km\,s$^{-1}$ from optical sodium \citep{wyttenbach2015} versus
${\sim}1.7$~km\,s$^{-1}$ from near-infrared CO/H$_2$O
\citep{brogi2016}. The modal relations are \textit{supports} and
\textit{complements} (e.g., terminator transmission versus dayside
emission water results for HD~209458~b probe different geometries and do
not conflict), followed by \textit{challenges-method} (a 1D-retrieval
bias result \citep{macdonald2020} limits the interpretation of a water
abundance \citep{kreidberg2015} without negating the detection). An
early binary confirm/reject protocol would have discarded most of these
relationships or, worse, upgraded them to contradictions.

\begin{table}[t]
  \centering
  \caption{Human adjudication of the 16 reviewed tension candidates.}
  \label{tab:adjudication}
  \begin{tabular}{lr}
    \toprule
    supports & 7 \\
    complements & 4 \\
    challenges-method & 2 \\
    qualifies & 1 \\
    explains-discrepancy & 1 \\
    contradicts (confirmed observational tension) & 1 \\
    \bottomrule
  \end{tabular}
\end{table}

\subsection{Question generation and two-stage ranking}
\label{sec:experiments:ranking}

From the prioritized signals (1~confirmed tension, 2~methodological
challenges, 1~independent qualification, 7~trusted single-dataset
conclusions), the generator produced 11 candidates; editorial curation
merged one near-duplicate pair into a question family with two
sub-questions and rewrote three questions (one had conflated statistical
significance with physical conditions; one presupposed the tension's
favored explanation; one was reframed for historical validation). The
Stage-A gate additionally caught the presupposition defect
independently. After curation, all 10 questions pass the gate.

Stage-B ranking places the wind-discrepancy question and the terminator-
heterogeneity family effectively tied at the top (scientific priority
8.13 each) with opposite execution profiles: the confirmed tension is
scientifically first-equal but operationally harder (execution priority
6.9---the needed simultaneous multi-band data may not exist), whereas
the heterogeneity family is reanalysis-ready against 2{,}643 archival
HST spectra of WASP-12 (execution 9.1). Significance caps hold the seven
P4 robustness questions to a 6.6--7.1 band, below both tension-driven
questions.

\subsection{Historical validation}
\label{sec:experiments:validation}

The validation corpus contains 1{,}891 unique 2021--2026 papers
(collected under its own manifest; a contamination guard keeps
\corpusA{} tables bit-identical before and after collection). Verdicts
for the 10 ranked questions are summarized in
Table~\ref{tab:validation}.

\begin{table}[t]
  \centering
  \caption{Historical validation: fate of the 10 ranked questions
  (generated from pre-2021 evidence) in the 2021--2026 literature.}
  \label{tab:validation}
  \begin{tabular}{lr}
    \toprule
    answered & 2 \\
    posed but open & 1 \\
    partially addressed & 7 \\
    not addressed & 0 \\
    \bottomrule
  \end{tabular}
\end{table}

Three outcomes deserve emphasis. First, \textbf{a premise refuted}: the
question of whether HD~209458~b's strongly subsolar terminator water
abundance \citep{macdonald2017} is atmospheric reality or retrieval
artifact was answered by 2025 reanalyses using independent retrieval
frameworks with improved systematics treatment, which find
solar-consistent water---the pipeline had isolated a published
conclusion that was in fact an artifact awaiting correction. Second,
\textbf{the top question is independently posed and open}: the
terminator-heterogeneity family is actively pursued by the community
without resolution. Third, \textbf{the reframed historical-validation
case behaves as designed}: JWST/NIRSpec PRISM programs are testing the
pre-launch TRAPPIST-1 CO$_2$ detectability prediction
\citep{lustigyaeger2019}, with stellar contamination \citep{rackham2018}
currently blocking a decisive verdict---the judge correctly reports the
premise as untested rather than refuted.

That zero questions were ignored by the subsequent literature is the
central quantitative result: questions surfaced systematically from
evidence tensions in pre-2021 knowledge are the questions the field
subsequently invested in.
